\chapter{Prérequis et installation}
\minitoc

\section{Prérequis}

\subsection{Logiciels}

\begin{table}[h]
\centering
\begin{tabular}{@{}p{4cm}p{3cm}p{6cm}@{}}
\toprule
\textbf{Élément} & \textbf{Version} & \textbf{Remarque} \\
\midrule
Python & 3.11 ou plus & Obligatoire. La plateforme emploie des
constructions introduites en 3.11. \\[6pt]
Navigateur & récent & Chrome, Firefox ou Edge. La console est une
application web. \\[6pt]
Docker & facultatif & Chemin de déploiement recommandé en production. \\[6pt]
Git & facultatif & Pour récupérer le dépôt~; une archive convient aussi. \\
\bottomrule
\end{tabular}
\caption{Logiciels requis}
\end{table}

Aucune base de données externe n'est à installer : la plateforme emploie
SQLite, contenu dans Python.

\subsection{Matériel}

Pour un usage de démonstration ou de recette, deux cœurs et deux gigaoctets de
mémoire suffisent. En exploitation, dimensionner selon le volume d'observations
reçues~; l'empreinte mémoire de la plateforme reste modeste, l'essentiel du
coût étant l'écriture au journal.

\subsection{Réseau}

La plateforme fonctionne \textbf{entièrement hors ligne}. Aucun appel sortant
n'est nécessaire à son fonctionnement nominal : c'est une exigence du Centre,
et elle a commandé le choix de toutes les bibliothèques employées
(chapitre~8).

\section{Installation depuis le dépôt}

\begin{lstlisting}[language=bash]
$ git clone https://github.com/Freddy-Gity/CIRTDEFENSE.git
$ cd CIRTDEFENSE
$ python3 -m venv .venv
$ source .venv/bin/activate
$ pip install -e ".[dev]"
\end{lstlisting}

Sous Windows, dans PowerShell :

\begin{lstlisting}[language=bash]
PS> py -3.11 -m venv .venv
PS> .\.venv\Scripts\Activate.ps1
PS> pip install -e ".[dev]"
\end{lstlisting}

\begin{tcolorbox}[colback=orange!5!white, colframe=orange!70!black,
  title=\textbf{Attention --- installez en mode éditable}, fonttitle=\bfseries,
  boxrule=0.6pt]
Le tiret \texttt{-e} n'est pas facultatif. L'interface web et l'emblème
imprimé en tête des rapports vivent dans le dossier \texttt{web/}, qui n'est
pas déclaré comme donnée du paquet. Une installation ordinaire
(\texttt{pip install .}) installe le code mais \textbf{pas} l'interface : le
serveur démarre, l'API répond, et la console est introuvable. Les guillemets
autour de \texttt{".[dev]"} sont nécessaires sous PowerShell, à cause des
crochets.
\end{tcolorbox}

\section{Vérification de l'installation}

Trois commandes suffisent à établir que tout est en place.

\begin{lstlisting}[language=bash]
$ pytest                      # doit afficher : 561 passed
$ ruff check src tests        # doit afficher : All checks passed!
$ cirtd status                # posture d'autonomie et etat du systeme
\end{lstlisting}

Si \texttt{pytest} passe, la chaîne complète — ingestion, décision, exécution,
annulation, journal, rapports — fonctionne sur votre poste.

\section{Démarrage}

\begin{lstlisting}[language=bash]
$ cirtd serve --host 127.0.0.1 --port 8000
\end{lstlisting}

Puis ouvrir \texttt{http://127.0.0.1:8000} dans un navigateur.

\begin{tcolorbox}[colback=green!4!white, colframe=darkgreen,
  title=\textbf{Bon à savoir}, fonttitle=\bfseries, boxrule=0.6pt]
Par défaut la plateforme écoute sur \texttt{0.0.0.0}, c'est-à-dire sur toutes
les interfaces. En poste de travail, préférez \texttt{127.0.0.1} : certaines
configurations Windows refusent l'écoute sur toutes les interfaces à un compte
non administrateur, et le message d'erreur n'est pas explicite
(voir chapitre~9).
\end{tcolorbox}

\section{Déploiement par conteneur}

\begin{lstlisting}[language=bash]
$ docker build -t cirtdefense .
$ docker run -p 8000:8000 -v cirt-data:/data cirtdefense
\end{lstlisting}

Le volume monté sur \texttt{/data} est \textbf{indispensable}. Sans lui, un
redémarrage perdrait le journal d'audit et relancerait l'autonomie que le
coupe-circuit venait d'interrompre.

\begin{tcolorbox}[colback=red!4!white, colframe=darkred,
  title=\textbf{Limitation connue}, fonttitle=\bfseries, boxrule=0.6pt]
Dans la topologie décrite par \texttt{docker-compose.yml}, la boucle de
contrôle qui mesure l'effet des gestes s'exécute dans un conteneur distinct de
celui qui les applique. Elle ne dispose donc pas des mesures de référence
prises au moment de l'action, et ne vérifie rien --- silencieusement. Tant que
ce point n'est pas corrigé, exécutez la boucle de contrôle dans le même
processus que l'API, ou lancez-la par la commande \texttt{cirtd control-loop}
sur le conteneur applicatif.
\end{tcolorbox}

\section{Configuration}

Toute la configuration passe par des variables d'environnement. Aucune n'est
obligatoire : la plateforme démarre avec des valeurs par défaut prudentes,
notamment le \emph{mode simulation}, dans lequel les gestes sont décidés et
tracés mais n'ont aucun effet réel sur les équipements.

Les trois variables à connaître avant tout :

\begin{table}[h]
\centering
\begin{tabular}{@{}p{5.4cm}p{2.4cm}p{5.6cm}@{}}
\toprule
\textbf{Variable} & \textbf{Défaut} & \textbf{Rôle} \\
\midrule
\texttt{CIRT\_ACTUATION\_MODE} & \texttt{simulation} &
\texttt{live} pour agir réellement sur les équipements. \\[6pt]
\texttt{CIRT\_AUTONOMY\_ENABLED} & \texttt{true} &
\texttt{false} : la plateforme raisonne et journalise, mais n'exécute rien. \\[6pt]
\texttt{CIRT\_DB\_PATH} & \texttt{data/cirtdefense.db} &
Emplacement de la base. \\
\bottomrule
\end{tabular}
\caption{Variables essentielles}
\end{table}

La liste complète des vingt-sept variables figure en annexe~A.

\begin{tcolorbox}[colback=orange!5!white, colframe=orange!70!black,
  title=\textbf{Attention --- passage en production}, fonttitle=\bfseries,
  boxrule=0.6pt]
Avant de basculer \texttt{CIRT\_ACTUATION\_MODE} sur \texttt{live}, changez
impérativement \texttt{CIRT\_ADMIN\_TOKEN} et \texttt{CIRT\_ANALYST\_TOKEN} :
leurs valeurs par défaut (\texttt{change-me-admin},
\texttt{change-me-analyst}) sont publiques. Un jeton par défaut en production
donne à quiconque le droit de déclencher des gestes sur vos équipements.
\end{tcolorbox}
